\part*{Chimie (7 points) : Suivi temporel d'une réaction d'oxydo-réduction \\
 Analyse d'un comprimé d'acide ascorbique}
\begin{center}
    \textbf{Les deux parties sont indépendantes}
\end{center}
Les réactions d'oxydo-réduction et d'acide-base sont deux types de
transformations chimiques d'une importance capitale en chimie des solutions.

Ces transformations peuvent être étudiées par différentes méthodes, ce qui
permet de faire un suivi temporel d'un système chimique, et déterminer certaines
caractéristiques et grandeurs\ldots Cet exercice vise~:
\begin{itemize}
    \item le suivi temporel d'une réaction d'oxydo-réduction~;
    \item l'analyse d'un comprimé d'acide ascorbique.
\end{itemize}

\section{Suivi temporel d'une réaction d'oxydo-réduction}
À l'instant $t_{0}=0$, on prépare une solution ($S$) en mélangeant un volume
d'une solution aqueuse d'iodure de potassium
($\ce{K^+_{(aq)} + I^-_{(aq)}}$)
contenant $n_{1}=\qty{8e-2}{\mol}$ d'ions $\ce{I^-_{(aq)}}$, avec
un volume d'une solution aqueuse de peroxodisulfate de sodium
($\ce{2Na^+_{(aq)} + S2O^2-_{8(aq)}}$)
contenant $n_{2}=\qty{2e-2}{\mol}$ d'ions $\ce{S2O^2-_{8(aq)}}$. Le volume
totale de la solution est $V=\qty{200}{\mL}$. Au cours de
la réaction, il se forme le diiode selon l'équation-bilan~:
\[
    \ce{S2O^2-_{8(aq)} + 2I^-_{(aq)} -> 2SO^2-_{4(aq)} + I2_{(aq)}}
\]

\begin{enumerate}
    \item Déterminer la valeur de l'avancement maximale $x_{\max}$. Déduire le
          réactif limitant.
    \item La courbe de la figure ci-dessous donne les variations de la quantité
          de matière du diiode formé en fonction du temps
          $n(\ce{I2})=f(t)$.
          \begin{figure}[H]
              \centering
              \includegraphics[width=0.7\textwidth]{2025_11_26_673124dc142c0df571a8g-2}
          \end{figure}
          
          \begin{enumerate}
              \item Calculer en (\unit{\mol\per\L\per\min}), la valeur de la
                    vitesse volumique de réaction à l'instant $t_{0}=0$.
              \item La valeur de la vitesse volumique de la réaction à l'instant
                    ${t_{1}=\qty{18}{\min}}$ est
                    ${v_{1}=\qty{1.44e-3}{\mol\per\L\per\min}}$.

                    Expliquer la diminution de la vitesse volumique de la
                    réaction.
              \item Citer un facteur cinétique permettant d'augmenter la vitesse
                    volumique de réaction sans changer l'état initial du système
                    chimique.
              \item Déterminer graphiquement le temps de demi-réaction $t_{1/2}$.
          \end{enumerate}
\end{enumerate}


\section{Analyse d'un comprimé d'acide ascorbique}
L'acide ascorbique $\ce{C6H8O6}$, couramment dénommé vitamine C, peut se trouver
en pharmacie sous forme de comprimés labélisés <<Vitamine C 500>>.

\subsection{Étude d'une solution aqueuse d'acide ascorbique}
Une solution aqueuse d'acide ascorbique $\mathrm{C}_{6} \mathrm{H}_{8} \mathrm{O}_{6(a q)}$, de concentration molaire $\mathrm{C}=4.10^{-3} \mathrm{~mol} . \mathrm{L}^{-1}$ et de volume $V=100 m L$ a un $p H=3,25$ à $25^{\circ} C$. L'acide ascorbique réagit avec l'eau selon l'équation chimique : $\mathrm{C}_{6} \mathrm{H}_{8} \mathrm{O}_{6(a q)}+\mathrm{H}_{2} \mathrm{O}_{(l)} \rightleftharpoons \mathrm{C}_{6} \mathrm{H}_{7} \mathrm{O}_{6(a q)}^{-}+\mathrm{H}_{3} \mathrm{O}_{(a q)}^{+}$\\
1.1. Identifier les deux couples acide/base mis en jeu.\\
1.2. Dresser le tableau d'avancement de la réaction en utilisant les grandeurs $C, V$, l'avancement $x$ et l'avancement $x_{\text {éq }}$ à l'état d'équilibre du système chimique.\\
1.3. Recopier sur votre copie le numéro de la question et écrire la lettre correspondante à la proposition vraie.\\
Le taux d'avancement final de la réaction vaut :

\begin{center}
\begin{tabular}{|l|c|l|l|l|l|l|l|}
\hline
$\mathbf{A}$ & $\tau \approx 0,34$ & $\mathbf{B}$ & $\tau \approx 0,47$ & $\mathbf{C}$ & $\tau \approx 0,55$ & $\mathbf{D}$ & $\tau \approx 0,14$ \\
\hline
\end{tabular}
\end{center}

1.4. Recopier sur votre copie le numéro de la question et écrire la lettre correspondante à la proposition vraie.\\
Le taux d'avancement final de la transformation dépend :

\begin{center}
\begin{tabular}{|c|l|}
\hline
A & \begin{tabular}{l}
de la constante d'équilibre $K$ associée à l'équation de la réaction et de la composition \\
initiale du système chimique \\
\end{tabular} \\
\hline
B & de la composition initiale du système chimique uniquement \\
\hline
C & de la constante d'équilibre $K$ associée à l'équation de la réaction uniquement \\
\hline
D & de la température du système chimique uniquement \\
\hline
\end{tabular}
\end{center}

1.5. Montrer que l'expression de la constante d'équilibre $K$ associée à l'équation de la réaction s'écrit : $K=\frac{C . \tau^{2}}{1-\tau}$. Calculer la constante d'acidité $K_{A}$ du couple $C_{6} H_{8} O_{6(a q)} / \mathrm{C}_{6} \mathrm{H}_{7} \mathrm{O}_{6(a q)}^{-}$.



\subsection{Vérification de la masse d'acide ascorbique dans un comprimé}
On écrase un comprimé de «Vitamine C 500», et on dissout la poudre dans l'eau pour obtenir une solution aqueuse ( $S_{A}$ ) de volume $V_{0}=200 m L$ et de concentration molaire $C_{A}$.\\
On dose le volume $V_{A}=20 m L$ de la solution ( $S_{A}$ ), avec une solution aqueuse d'hydroxyde de sodium $N a_{(a q)}^{+}+H O_{(a q)}^{-}$de concentration molaire $C_{B}=2,0.10^{-2} m o l . L^{-1}$. L'équivalence est obtenue pour un volume de solution versée $V_{B, E}=14,2 m L$.\\
2.1. Écrire l'équation chimique modélisant la transformation qui a eu lieu lors du dosage.\\
2.2. Calculer la concentration molaire $C_{A}$.\\
2.3. En déduire la valeur de la masse d'acide ascorbique contenu dans ce comprimé puis expliquer l'indication «Vitamine C 500». On donne : $M\left(C_{6} H_{8} O_{6}\right)=176 g$. mol $^{-1}$.



